% GENERATED FILE. Edit canonical lessons, then run npm run generate:books.
% canonical-source-hash: fnv1a64:8083ff5e033e622d
% canonical-lessons: TA-C302-peralivu, TA-C302-uyirilappu, TA-C302-nivaranam, TA-C302-mintelu, TA-C302-curaiyatu

\chapter{After the Disaster}
\label{ch:ta-after-the-disaster}
\hlchaptermodality{302}

\begin{chapteropening}
\textbf{By the end of this chapter:} \emph{I can talk about a disaster and what came after it.}

Complete the last lesson of chapter 302: I can talk about a disaster and what came after it.
\end{chapteropening}

\section[\texorpdfstring{\ta{பேரழிவு} (pēraḻivu)}{பேரழிவு (pēraḻivu)}]{\ta{பேரழிவு} (pēraḻivu) — a disaster}

\label{lesson:TA-C302-peralivu}



\begin{quote}
Before the new one: say the Tamil for to survive, then the Tamil for to retrieve.
\end{quote}

\subsection*{You'll want to know: \ta{பேரழிவு}}
\textbf{\ta{பேரழிவு}} — \emph{pēraḻivu} — \textquotedblleft{}a disaster\textquotedblright{}.

\begin{quote}
\ta{வெள்ளம்} \ta{ஒரு} \ta{பேரழிவாக} \ta{மாறியது}.
\end{quote}

\begin{quote}
\emph{veḷḷam oru pēraḻivāka māṟiyatu.}
\end{quote}

\begin{quote}
The flood turned into a disaster.
\end{quote}

\ta{பேர்}- before a noun means great.

\begin{cousinweb}
\textbf{\ta{பேரழிவு}} (\emph{pēraḻivu}) is a Tamil compound built from \ta{பேர்}, \textquotedblleft{}great\textquotedblright{} (DEDR 4411), and \ta{அழிவு}, \textquotedblleft{}destruction, ruin\textquotedblright{} (DEDR 277).
\end{cousinweb}

\begin{grammarlens}[title={What you've built}]
One more word for telling what happened.
\end{grammarlens}

\subsection*{Your turn}
\begin{itemize}
\raggedright
  \item \emph{Say it:} \emph{pēraḻivu}
  \item \emph{Say it:} \emph{pēraḻivu}, once more
  \item \emph{Say it:} \emph{veḷḷam oru pēraḻivāka māṟiyatu.} — The flood turned into a disaster.
  \item \emph{Recall:} say the Tamil for to survive, then the Tamil for to retrieve, then say \emph{pēraḻivu} again
\end{itemize}

\begin{tcolorbox}[breakable,colback=teal!4,colframe=teal!35!black,title={Before you move on}]
What does \ta{பேரழிவு} mean? (\textquotedblleft{}A disaster\textquotedblright{}.) Say it once more.
\end{tcolorbox}
\section[\texorpdfstring{\ta{உயிரிழப்பு} (uyiriḻappu)}{உயிரிழப்பு (uyiriḻappu)}]{\ta{உயிரிழப்பு} (uyiriḻappu) — a loss of life}

\label{lesson:TA-C302-uyirilappu}



\begin{quote}
Before the new one: say the Tamil for to retrieve, then the Tamil for a disaster.
\end{quote}

\subsection*{You'll want to know: \ta{உயிரிழப்பு}}
\textbf{\ta{உயிரிழப்பு}} — \emph{uyiriḻappu} — \textquotedblleft{}a loss of life\textquotedblright{}.

\begin{quote}
\ta{நல்லவேளையாக}, \ta{உயிரிழப்பு} \ta{எதுவும்} \ta{இல்லை}.
\end{quote}

\begin{quote}
\emph{nallavēḷaiyāka, uyiriḻappu etuvum illai.}
\end{quote}

\begin{quote}
Fortunately, there was no loss of life.
\end{quote}

\ta{உயிர்}, life + \ta{இழப்பு}, loss, from \ta{இழ}, to lose.

\begin{cousinweb}
\textbf{\ta{உயிரிழப்பு}} (\emph{uyiriḻappu}) is a Tamil compound built from \ta{உயிர்}, \textquotedblleft{}life\textquotedblright{} (DEDR 645), and \ta{இழப்பு}, \textquotedblleft{}loss\textquotedblright{}, from \ta{இழ}, \textquotedblleft{}to lose\textquotedblright{} (DEDR 501).
\end{cousinweb}

\begin{grammarlens}[title={What you've built}]
One more word for telling what happened.
\end{grammarlens}

\subsection*{Your turn}
\begin{itemize}
\raggedright
  \item \emph{Say it:} \emph{uyiriḻappu}
  \item \emph{Say it:} \emph{uyiriḻappu}, once more
  \item \emph{Say it:} \emph{nallavēḷaiyāka, uyiriḻappu etuvum illai.} — Fortunately, there was no loss of life.
  \item \emph{Recall:} say the Tamil for to retrieve, then the Tamil for a disaster, then say \emph{uyiriḻappu} again
\end{itemize}

\begin{tcolorbox}[breakable,colback=teal!4,colframe=teal!35!black,title={Before you move on}]
What does \ta{உயிரிழப்பு} mean? (\textquotedblleft{}A loss of life\textquotedblright{}.) Say it once more.
\end{tcolorbox}
\section[\texorpdfstring{\ta{நிவாரணம்} (nivāraṇam)}{நிவாரணம் (nivāraṇam)}]{\ta{நிவாரணம்} (nivāraṇam) — relief aid; relief}

\label{lesson:TA-C302-nivaranam}



\begin{quote}
Before the new one: say the Tamil for a disaster, then the Tamil for a loss of life.
\end{quote}

\subsection*{You'll want to know: \ta{நிவாரணம்}}
\textbf{\ta{நிவாரணம்}} — \emph{nivāraṇam} — \textquotedblleft{}relief aid; relief\textquotedblright{}.

\begin{quote}
\ta{அரசு} \ta{மக்களுக்கு} \ta{நிவாரணம்} \ta{கொடுத்தது}.
\end{quote}

\begin{quote}
\emph{aracu makkaḷukku nivāraṇam koṭuttatu.}
\end{quote}

\begin{quote}
The government gave the people relief aid.
\end{quote}

Also relief from pain: \ta{வலிக்கு} \ta{நிவாரணம்}.

\begin{cousinweb}
\textbf{\ta{நிவாரணம்}} (\emph{nivāraṇam}) is borrowed from Sanskrit \emph{nivāraṇa}, \textquotedblleft{}keeping back, preventing\textquotedblright{} (CDIAL 7418): what keeps harm away.
\end{cousinweb}

\begin{grammarlens}[title={What you've built}]
One more word for telling what happened.
\end{grammarlens}

\subsection*{Your turn}
\begin{itemize}
\raggedright
  \item \emph{Say it:} \emph{nivāraṇam}
  \item \emph{Say it:} \emph{nivāraṇam}, once more
  \item \emph{Say it:} \emph{aracu makkaḷukku nivāraṇam koṭuttatu.} — The government gave the people relief aid.
  \item \emph{Recall:} say the Tamil for a disaster, then the Tamil for a loss of life, then say \emph{nivāraṇam} again
\end{itemize}

\begin{tcolorbox}[breakable,colback=teal!4,colframe=teal!35!black,title={Before you move on}]
What does \ta{நிவாரணம்} mean? (\textquotedblleft{}Relief aid; relief\textquotedblright{}.) Say it once more.
\end{tcolorbox}
\section[\texorpdfstring{\ta{மீண்டெழு} (mīṇṭeḻu)}{மீண்டெழு (mīṇṭeḻu)}]{\ta{மீண்டெழு} (mīṇṭeḻu) — to rise again}

\label{lesson:TA-C302-mintelu}



\begin{quote}
Before the new one: say the Tamil for a loss of life, then the Tamil for relief aid.
\end{quote}

\subsection*{You'll want to know: \ta{மீண்டெழு}}
\textbf{\ta{மீண்டெழு}} — \emph{mīṇṭeḻu} — \textquotedblleft{}to rise again\textquotedblright{}.

\begin{quote}
\ta{சில} \ta{மாதங்களில்} \ta{அந்த} \ta{ஊர்} \ta{மீண்டெழுந்தது}.
\end{quote}

\begin{quote}
\emph{cila mātaṅkaḷil anta ūr mīṇṭeḻuntatu.}
\end{quote}

\begin{quote}
In a few months that town rose again.
\end{quote}

Past (I …): \ta{மீண்டெழுந்தேன்}, \emph{mīṇṭeḻuntēṉ}.

From \ta{மீண்டு}, having returned, and \ta{எழு}, to rise.

\begin{cousinweb}
\textbf{\ta{மீண்டெழு}} (\emph{mīṇṭeḻu}) is a Tamil compound built from \ta{மீண்டு}, \textquotedblleft{}having returned\textquotedblright{}, from \ta{மீள்} (DEDR 4883), and \ta{எழு}, \textquotedblleft{}to rise\textquotedblright{} (DEDR 851).
\end{cousinweb}

\begin{grammarlens}[title={What you've built}]
One more word for telling what happened.
\end{grammarlens}

\subsection*{Your turn}
\begin{itemize}
\raggedright
  \item \emph{Say it:} \emph{mīṇṭeḻu}
  \item \emph{Say it:} \emph{mīṇṭeḻu}, once more
  \item \emph{Say it:} \emph{cila mātaṅkaḷil anta ūr mīṇṭeḻuntatu.} — In a few months that town rose again.
  \item \emph{Recall:} say the Tamil for a loss of life, then the Tamil for relief aid, then say \emph{mīṇṭeḻu} again
\end{itemize}

\begin{tcolorbox}[breakable,colback=teal!4,colframe=teal!35!black,title={Before you move on}]
What does \ta{மீண்டெழு} mean? (\textquotedblleft{}To rise again\textquotedblright{}.) Say it once more.
\end{tcolorbox}
\section[\texorpdfstring{\ta{சூறையாடு} (cūṟaiyāṭu)}{சூறையாடு (cūṟaiyāṭu)}]{\ta{சூறையாடு} (cūṟaiyāṭu) — to ravage; to plunder}

\label{lesson:TA-C302-curaiyatu}



\begin{quote}
Before the new one: say the Tamil for relief aid, then the Tamil for to rise again.
\end{quote}

\subsection*{You'll want to know: \ta{சூறையாடு}}
\textbf{\ta{சூறையாடு}} — \emph{cūṟaiyāṭu} — \textquotedblleft{}to ravage; to plunder\textquotedblright{}.

\begin{quote}
\ta{வெள்ளம்} \ta{பல} \ta{கிராமங்களைச்} \ta{சூறையாடியது}.
\end{quote}

\begin{quote}
\emph{veḷḷam pala kirāmaṅkaḷaic cūṟaiyāṭiyatu.}
\end{quote}

\begin{quote}
The flood ravaged many villages.
\end{quote}

Past (I …): \ta{சூறையாடினேன்}, \emph{cūṟaiyāṭiṉēṉ}.

What is ravaged takes -\ta{ஐ}; said of robbers too.

\begin{cousinweb}
\textbf{\ta{சூறையாடு}} (\emph{cūṟaiyāṭu}) is a Tamil compound built from \ta{சூறை}, \textquotedblleft{}robbery, pillage\textquotedblright{} (DEDR 2744), and \ta{ஆடு}, \textquotedblleft{}to move, act\textquotedblright{} (DEDR 347).
\end{cousinweb}

\begin{grammarlens}[title={What you've built}]
One more word for telling what happened.
\end{grammarlens}

\subsection*{Your turn}
\begin{itemize}
\raggedright
  \item \emph{Say it:} \emph{cūṟaiyāṭu}
  \item \emph{Say it:} \emph{cūṟaiyāṭu}, once more
  \item \emph{Say it:} \emph{veḷḷam pala kirāmaṅkaḷaic cūṟaiyāṭiyatu.} — The flood ravaged many villages.
  \item \emph{Recall:} say the Tamil for relief aid, then the Tamil for to rise again, then say \emph{cūṟaiyāṭu} again
\end{itemize}

\begin{tcolorbox}[breakable,colback=teal!4,colframe=teal!35!black,title={Before you move on}]
What does \ta{சூறையாடு} mean? (\textquotedblleft{}To ravage; to plunder\textquotedblright{}.) Say it once more.
\end{tcolorbox}
